% Test-case specification TC-0002 (CTest test TC-0002-helpers-memory).
\documentclass[letterpaper]{article}

\usepackage[letterpaper,margin=1in]{geometry}
\emergencystretch=3em

% Compile from the repository root so this path resolves to
% the pinned WSP submodule.
\input{wsp/testing/test-case-library.tex}

\begin{document}

\def\TCID{TC-0002}
\def\TCTitle{%
  Member and memory helpers by name}
\def\TCPurpose{%
  Verify the member and memory helpers called by name from the built DLL:
  \texttt{GetMem1}/\texttt{2}/\texttt{4}/\texttt{8},
  \texttt{PutMem1}/\texttt{2}/\texttt{4}/\texttt{8}, \texttt{VarPtr}, the
  String, object and Variant forms.}
\def\TCPriority{\TCPriorityHigh}

\def\TCRequirementRef{%
  REQ-0005}
\def\TCDesignRef{%
  Not applicable.}
\def\TCFeatureRef{%
  \texttt{tests/\allowbreak{}vbrun\_\allowbreak{}helper\_\allowbreak{}tests.\allowbreak{}cpp}
  (family \texttt{memory}),
  \texttt{src/\allowbreak{}helpers\_\allowbreak{}memory.\allowbreak{}cpp}}
\def\TCTestTechnique{%
  Requirements-based testing with expected values probed on the reference
  runtime (\texttt{scripts/\allowbreak{}Probe-Vbrun.\allowbreak{}ps1}).}

\def\TCPreconditions{%
  The library and the test program are built
  (\texttt{scripts/\allowbreak{}Invoke-Build.\allowbreak{}ps1}) and the test
  is registered with CTest.}
\def\TCEnvironment{%
  Windows with the toolchain of
  \texttt{scripts/\allowbreak{}Invoke-Build.\allowbreak{}ps1}; the local
  gate builds and runs x64 Debug, CI builds and runs the other
  architectures. The library is built from this repository in
  \texttt{out/\allowbreak{}build}.}
\def\TCAssumptions{%
  Not applicable.}

\def\TCInputData{%
  Byte, word, long and eight-byte values such as \texttt{0xA5},
  \texttt{0xBEEF}, \texttt{0xDEADBEEF}; the text \texttt{hello}, whose odd
  byte length must be kept.}
\def\TCInitialState{%
  No state from earlier runs is needed; the test makes what it needs.}

\def\TCProcedure{%
  \begin{enumerate}
    \item CTest runs \texttt{vbrun\_\allowbreak{}helper\_\allowbreak{}tests
      \textless{}wfcvbrun.\allowbreak{}dll\textgreater{} memory}.
    \item Load the DLL and look up each export by name.
    \item Read values with the \texttt{GetMem} helpers and write them with the
      \texttt{PutMem} helpers and compare them with the source values.
    \item Check that \texttt{VarPtr} returns the address it was given.
    \item Copy Strings, objects and Variants with the member helpers and check
      the values, byte lengths and reference counts.
    \item CTest records the exit code and, on failure, the program's output
      (\texttt{--output-on-failure}).
  \end{enumerate}}
\def\TCExpectedResult{%
  Every check passes and the program prints \texttt{memory ok: N checks} and
  exits 0; a failed check prints its name and the exit code is 1.}
\def\TCPostConditions{%
  No file or registry state remains that later tests depend on.}

\TCTemplate

\end{document}
